\documentclass[11pt]{article}
\usepackage[a4paper,margin=27mm]{geometry}
\usepackage{amsmath,amssymb,mathtools}
\usepackage{fontspec,unicode-math}
\setmainfont{Latin Modern Roman}
\setmathfont{Latin Modern Math}
\usepackage{xurl}
\usepackage[hidelinks]{hyperref}
\newcommand{\sref}[2]{\hyperlink{source:#1:#2}{[S#2]}}
\newcommand{\cataloguescope}[1]{\paragraph{Mathematical question.}#1}
\hypersetup{pdftitle={UnsolvedMath Problem 9400170}}
\begin{document}
% UnsolvedMath ID 9400170; source catalogue code AMR-093-0170
\setcounter{section}{9400169}
\section{What is the smallest Skewes's number}\label{9400170}
{\small\sffamily UnsolvedMath ID 9400170 · Number Theory}
\subsection{Definitions and mathematical statement}

\paragraph{Shared notation.}$\mathbb N=\{1,2,\ldots\}$, and $\mathbb Z,\mathbb Q,\mathbb R,\mathbb C$ denote the integers, rationals, reals and complex numbers. Any other conventions are specified locally.


For a real $x\ge2$, let $\pi(x)=\#\{p\le x:p\text{ is a positive prime}\}$ and define the unshifted logarithmic integral by
\[\operatorname{li}(x)=\lim_{\delta\downarrow0}\left(\int_0^{1-\delta}\frac{dt}{\log t}+\int_{1+\delta}^{x}\frac{dt}{\log t}\right).\]
Thus $\operatorname{li}$ uses Cauchy principal value at $1$, rather than the offset integral starting at $2$. The least natural-number crossover is
\[x_0=\min\{n\in\mathbb N:n\ge2,\ \pi(n)>\operatorname{li}(n)\}.\]
Littlewood's sign-change theorem guarantees that the displayed set is nonempty.
\paragraph{Statement recorded in the supplied source.}
What is the exact value of $x_0$, the smallest natural number at which $\pi(x)$ exceeds $\operatorname{li}(x)$? \sref{9400170}{2}
\subsection{Short English statement}
Determine the very first integer where the prime-counting function exceeds the unshifted logarithmic integral.
\subsection{Sources}
\begin{description}
\item[\hypertarget{source:9400170:1}{S1}] ULAM.AI and the UnsolvedMath Contributors, UnsolvedMath: A Curated Collection of Open Mathematics Problems, record 9400170 (AMR-093-0170). CC BY 4.0. \url{https://huggingface.co/datasets/ulamai/UnsolvedMath}.
\item[\hypertarget{source:9400170:2}{S2}] Wikipedia contributors, \emph{Skewes's number}, revision 1370942184 (2026) Definition and evidence. \url{https://en.wikipedia.org/w/index.php?title=Skewes%27s+number&oldid=1370942184}.
\item[\hypertarget{source:9400170:3}{S3}] Wikipedia contributors, \emph{Logarithmic integral function}, revision 1339812367 (2026) Integral representation; Offset logarithmic integral. \url{https://en.wikipedia.org/w/index.php?title=Logarithmic+integral+function&oldid=1339812367}.
\item[\hypertarget{source:9400170:4}{S4}] Wikipedia contributors, \emph{List of unsolved problems in mathematics}, revision 1366636767 (2026) Number theory, What is the smallest Skewes's number. \url{https://en.wikipedia.org/w/index.php?title=List_of_unsolved_problems_in_mathematics&oldid=1366636767#Number_theory}.
\end{description}
\label{end:9400170}
\end{document}
